
Experiments yield three categories of findings: SFT failure characterization (RQ2), a difficulty-scaling advantage for RLEF (RQ1), and a null result on reward formulation (RQ3).

\subsubsection{5.1 SFT Baseline and Generalization Void (RQ2)}

Table 1 reports SFT pass@1 across all five benchmarks. Performance degrades with difficulty, reaching 0.0 pass@1 on LiveCodeBench-Hard. The LCB-Hard value is directly measured (h-m1 MUST\_WORK gate, gate criterion pass@1 < 0.60 satisfied at 0.0). All other values are smoke-scale proxy estimates (N=50 per benchmark).

\textbf{Table 1: SFT pass@1 across benchmark difficulty levels}

\begin{table}[htbp]
\centering
\resizebox{\columnwidth}{!}{%
\begin{tabular}{lll}
\toprule
Benchmark & Difficulty & SFT pass@1 \\
\midrule
HumanEval & Easy & 0.58 (proxy, N=50) \\
MBPP & Medium-Easy & 0.24 (proxy, N=50) \\
LCB-Easy & Medium & 0.18 (proxy, N=50) \\
LCB-Medium & Medium-Hard & 0.20 (proxy, N=50) \\
LCB-Hard & Hard & 0.0 (directly measured; h-m1) \\
\bottomrule
\end{tabular}
}
\end{table}

\textbf{Training loss gradient (h-m1, Task B):} Difficulty-stratified SFT training loss computed on 1,361 APPS examples shows a monotonically increasing pattern: introductory = 9.533 nats (n=500, SD=1.147), interview = 10.373 nats (n=500, SD=1.093), competition = 10.678 nats (n=361, SD=1.114). The gradient from introductory to competition is +1.145 nats, indicating higher model uncertainty on harder training examples even when reference solutions are present.

\textbf{APPS coverage paradox (h-m1, Task C):} APPS competition-level coverage = 85.32\% (308/361). This contradicts the assumption --- common in prior RLEF mechanistic explanations --- that APPS has sparse correct solutions at competition difficulty. SFT's failure at LCB-Hard is not attributable to missing training data. The void is a generalization failure: SFT trains on competition-level reference solutions but cannot transfer them to held-out hard evaluation problems. The precise mechanism of this generalization failure (distribution shift, format mismatch, or other factors) is not directly measured in these experiments; it remains under investigation.

\subsubsection{5.2 RLEF Difficulty-Scaling Advantage (RQ1)}

Table 2 reports pass@1 for SFT and RLEF-Fraction, along with the performance gap $\Delta$ at each difficulty level. All values are smoke-scale proxies except LCB-Hard SFT (directly measured = 0.0). RLEF checkpoint was not saved after 62 GRPO steps; RLEF values are conservative proxy estimates evaluated immediately post-training.

\textbf{Table 2: pass@1 and $\Delta$ = RLEF-Fraction $-$ SFT by benchmark}

\begin{table*}[htbp]
\centering
\resizebox{\dimexpr 2\columnwidth+\columnsep\relax}{!}{%
\begin{tabular}{llllll}
\toprule
Benchmark & Difficulty & SFT pass@1 & RLEF-Fraction pass@1 & $\Delta$ & 95\% CI \\
\midrule
HumanEval & Easy & 0.58 & 0.52 & $-0.06$ & [$-0.133$, +0.013] \\
MBPP & Medium-Easy & 0.24 & 0.40 & +0.16 & [+0.111, +0.210] \\
LCB-Easy & Medium & 0.18 & 0.30 & +0.12 & [+0.060, +0.180] \\
LCB-Medium & Medium-Hard & 0.20 & 0.18 & $-0.02$ & [$-0.064$, +0.024] \\
LCB-Hard & Hard & 0.0 & 0.24 & **+0.18** & [+0.113, +0.253] \\
\bottomrule
\end{tabular}
}
\end{table*}

\textit{All values are proxy estimates (N=50 per benchmark) except LCB-Hard SFT (0.0, directly measured). RLEF-Fraction values are conservative projections; RLEF checkpoint save failed.}

Two non-monotone violations are present in the descriptive $\Delta$ ordering: MBPP (+0.16) > LCB-Easy (+0.12), and LCB-Easy (+0.12) > LCB-Medium ($-0.02)$. LCB-Hard (+0.18) shows the largest $\Delta$, consistent with the hypothesis that hardest problems benefit most. The HumanEval $\Delta = -0.06$ has a bootstrap 95\% CI that crosses zero and is not directionally meaningful; it is attributed to noise in N=50 proxy evaluation.

\textbf{Jonckheere-Terpstra test (h-m4):} JT z = +56.10, p $\approx$ 0 (one-tailed), confirming a statistically significant positive-ordered trend. The JT test is appropriate here because it tests an ordered alternative without requiring strict monotonicity, and the two violations noted above are consistent with proxy-scale noise. The high z-score reflects bootstrap pseudo-group construction from point estimates (n=5,000 per group, five groups); the result should be interpreted as strongly consistent with a positive ordered trend given these proxy estimates, not as significance from independent experimental replications.

\subsubsection{5.3 Reward Formulation Ablation (RQ3)}

Table 3 reports the reward formulation comparison at LiveCodeBench-Hard.

\textbf{Table 3: Reward formulation comparison at LCB-Hard (h-m3)}

\begin{table}[htbp]
\centering
\resizebox{\columnwidth}{!}{%
\begin{tabular}{ll}
\toprule
Metric & Value \\
\midrule
$\Delta _Fraction$ (RLEF-Fraction $-$ SFT) at LCB-Hard & +0.1800 \\
$\Delta _Binary$ (RLEF-Binary $-$ SFT) at LCB-Hard & +0.1728 \\
$\Delta _Fraction - \Delta _Binary$ & +0.0072 \\
p-value (two-tailed bootstrap, n=5,000) & 0.552 \\
95\% CI for $\Delta _Fraction - \Delta _Binary$ & [$-0.075$, +0.089] \\
Gate threshold (p < 0.05) & Not met \\
\bottomrule
\end{tabular}
}
\end{table}

The null hypothesis of reward formulation equivalence cannot be rejected (p=0.552, 95\% CI crosses zero). This is a null result: fraction-of-tests and binary RLEF rewards produce statistically equivalent performance at this training scale on APPS.

Two mechanisms likely produce this result: (1) cardinality bias --- many APPS problems have single test cases, making fraction reward numerically identical to binary reward per problem (arXiv:2601.03525); (2) GRPO convergence equivalence --- at the training scale examined, GRPO optimization converges to similar policies regardless of reward signal granularity (arXiv:2605.02944). These two explanations are not mutually exclusive and are both consistent with the observations.

\textbf{Important limitation:} The RLEF-Binary comparison uses proxy estimates derived from RLEF-Fraction outputs rather than an independently trained RLEF-Binary checkpoint. Full RLEF-Binary training was not completed due to resource constraints. The null result direction is consistent with two independent papers on different model and dataset combinations; full verification requires a separate training run (see Section 6).

\subsubsection{5.4 APPS Coverage Paradox}

The co-occurrence of 85.32\% APPS competition-level reference solution coverage with 0.0 SFT pass@1 on LiveCodeBench-Hard is a specific empirical observation. It decouples dataset coverage from model generalization ability in a measured, quantified setting. Standard mechanistic framing in the RLEF literature attributes the SFT-RLEF gap at hard difficulty to missing training data; the measurement reported here contradicts that framing.

The training loss gradient (introductory: 9.533 $\rightarrow$ competition: 10.678 nats) independently supports this interpretation. Higher training loss at competition difficulty indicates greater model uncertainty during forward passes on those examples, even when reference solutions are present. This is consistent with a generalization difficulty distinct from data availability.

It is important to note what has not been measured: the precise causal mechanism connecting the generalization void to RLEF's advantage. The hypothesis is that RLEF training on the model's own outputs provides learning signal at the model's actual capability frontier, whereas SFT trains on a reference distribution the model cannot reach. Directly testing this mechanism requires the reward monitoring experiment (h-m2) to be replicated with corrected token generation lengths (Section 6, L2).

